\chapter{Code Availability and Software Provenance}
\label{ch:provenance}

The PPP guideline requires an explicit account of existing code, external
libraries, interfaces, and original contributions. The project uses mature
libraries for tensor differentiation and finite-element infrastructure, but
the constitutive bridge, homogenization workflow, rigid-particle treatment,
condensation logic, dataset generator, and verification suite were developed
for this project.

The NGSolve TensorFlow tutorial is a methodological precedent: it illustrates
how a differentiable constitutive calculation can be connected to finite
elements. It does not provide the PyTorch bridge used here, and TensorFlow is
not part of the reported runtime. The Phase~1 bridge evaluates the pICNN in
PyTorch at NGSolve integration points, transfers its energy derivatives to
NGSolve, and assembles the total-Lagrangian residual and tangent. In Phase~2,
the matrix and homogenization equations are instead expressed as NGSolve
symbolic forms and differentiated with NGSolve. These two derivative paths
serve distinct roles and are tested separately.

\section{Closest available resources}

\begin{table}[H]
\centering
\caption{Existing resources used as points of departure.}
\label{tab:available-code}
\begin{tabularx}{\textwidth}{@{}p{0.22\textwidth}X X@{}}
\toprule
Resource & Available functionality & Relationship to this project\\
\midrule
NGSolve tutorials and API
& Geometry, finite-element spaces, symbolic forms, periodic spaces, and
linear algebra
& Used as infrastructure. The periodic-space tutorial and TensorFlow example
provide relevant precedents, but the PyTorch material bridge, rigid RVE
condensation, and verification workflow are implemented here
\cite{ngsolve_periodic,ngsolve_tensorflow}.\\
PyTorch documentation
& Neural modules, optimizers, tensor operations, and automatic differentiation
& Used as infrastructure for the custom energy architecture and exact
stress/tangent evaluation; the coupling to NGSolve is project-specific.\\
ICNN literature and implementations
& Convex neural architectures and non-negative weight constraints
& Adapted to finite-strain invariant features, context conditioning, and
stress-free normalization.\\
Published neural hyperelastic models
& Physics-structured constitutive learning
& Used for methodological comparison; the NGSolve integration and RVE
pipeline are project-specific \cite{kalina2022,linden2023}.\\
\bottomrule
\end{tabularx}
\end{table}

\section{Original implementation}

The main original software components developed for this project are:
\begin{itemize}
  \item a PyTorch--NGSolve integration-point material bridge;
  \item exact first-Piola stress and full \(9\times9\) tangent extraction;
  \item periodic cell and particle-cavity geometry generation;
  \item mixed finite-strain matrix and custom Newton solution;
  \item exact finite rigid-particle kinematics and free rotation;
  \item reaction-based Hill--Mandel homogenization;
  \item energy-Hessian condensation of microscopic and rotational variables;
  \item deterministic, checkpointed RVE dataset generation; and
  \item executable Phase~1 and Phase~2 verification gates.
\end{itemize}

The constitutive model combines established ideas---invariant-based
hyperelasticity, input-convex networks, and automatic differentiation---but
the interface contract and FE material transfer are specific to this
implementation. Similarly, the periodic-space concept and Taylor--Hood
mixed discretization are established numerical tools; the project-specific
work is their coupling to an exactly rigid particle, free finite rotation,
reaction-based effective stress, a consistent condensed tangent, and a
restartable supervised dataset. The implementation does not claim that the
network architecture, finite-element library, or underlying mechanics
theory is original.

The source of truth for the numerical results is the code and its archived
configuration, not the initial presentation. Executable verification cases
are kept under \texttt{tests/}; the core equations and their FE assembly are
implemented in \texttt{src/}. NGSolve and PyTorch supply reusable library
functionality, while the data conversion, tensor-index convention, nonlinear
control flow, homogenization definitions, and error checks are explicit
project code.

\section{External libraries}

\begin{table}[H]
\centering
\caption{Principal external software and the version evidence retained for
the project environments.}
\label{tab:software}
\begin{tabular}{@{}lll@{}}
\toprule
Software & Role & Final version\\
\midrule
Python & Primary implementation language & 3.11.7\\
PyTorch & Neural model and automatic differentiation & 2.15.0+cpu\\
NGSolve/Netgen & FEM, geometry, symbolic differentiation & 6.2.2608\\
NumPy & Array storage and tensor post-processing & 2.4.6\\
Matplotlib & Verification plotting & 3.11.0 (Phase~1 baseline only)\\
\bottomrule
\end{tabular}
\end{table}

The Python, PyTorch, NGSolve, and NumPy versions are taken from the Phase~7
PBS scripts and captured standard output. Matplotlib 3.11.0 appears in the
Phase~1 baseline package, but a Phase~7 plotting-library version was not
recorded. An earlier local environment record contains Python~3.12.2,
PyTorch~2.5.1, NumPy~1.26.4, and NGSolve~6.2.2607; it is not used to label the
later cluster artifacts.
